\section{The claim and its setting}\label{sec:introduction}

Resolution over linear equations was introduced by Raz and Tzameret
\cite{RT08}, with equations over the integers. Its characteristic-two
version, resolution over parities, was subsequently studied by Itsykson
and Sokolov \cite{IS14,IS20}, who established tree-like lower bounds.
Here clauses are disjunctions of affine equations over \(\F\), with the
rules specified below. Its unrestricted proof-size lower-bound problem remains an
explicit open benchmark in the literature reviewed for this preprint
\cite{IPS26}. The contribution claimed here is an exponential lower
bound for the standard bit pigeonhole principle in unrestricted DAG-like
\(\Res\).

\subsection{The bit pigeonhole principle}
Fix \(\ell\ge2\), put \(n=2^\ell\) and \(m=n+1\), and introduce
\[
  b_{it}\qquad (i\in[m],\ t\in[\ell]).
\]
The row \(b_i\) encodes one of the \(n\) labels in \(\{0,1\}^{\ell}\).
For a label \(z\), let \([b_{it}\ne z_t]\) denote \(b_{it}\) when \(z_t=0\)
and \(\neg b_{it}\) when \(z_t=1\). The usual CNF encoding is
\begin{equation}\label{eq:bit-php}
 \BPHP_n^{n+1}
  =\bigwedge_{\substack{i<i'\in[m]\\z\in\{0,1\}^{\ell}}}
    \left(\bigvee_{t\in[\ell]}[b_{it}\ne z_t]
       \ \vee\ \bigvee_{t\in[\ell]}[b_{i't}\ne z_t]\right).
\end{equation}
It has \(v=m\ell\) variables and \(n\binom m2\) clauses of width \(2\ell\).
No separate range or pigeon axiom is required: every Boolean row encodes
exactly one available label.

\subsection{The proof system}
A \emph{linear clause} is a disjunction of affine equations over \(\F\).
The basic rules are complementary-parity resolution and semantic weakening:
\begin{equation}\label{eq:res-rules}
 \frac{A\vee(u=0)\qquad B\vee(u=1)}{A\vee B},
 \qquad
 \frac{C}{D}\quad\text{if } C\models D.
\end{equation}
The pivot \(u\) may be any affine form. A refutation is a finite acyclic
derivation from initial clauses to the empty clause. Its size is the number
of nodes, including initial nodes, in its proof DAG. Earlier nodes may be
used repeatedly, and no restriction is placed on proof depth or on repeated
use of related pivots.

We also consider the convention permitting every sound inference from at
most two linear-clause premises. Lemma~\ref{lem:affine-cover} reduces this
convention to \eqref{eq:res-rules} with constant node overhead.
Fresh unconstrained propositional variables, if used, may be set to
constants throughout a proof. Extension axioms defining new predicates
are not part of the system considered here.

\begin{theorem}[Main theorem]\label{thm:main}\lean{claims/BitPHPExponential.lean}
For every \(\ell\ge32\), with \(n=2^\ell\), every DAG-like \(\Res\)
refutation of \(\BPHP_{n}^{n+1}\) has more than
\[
 \exp\!\left(\frac{n}{32768\,\ell^2}\right)
\]
nodes. The same bound holds with base two in place of \(e\), and for both
rule conventions above.
\end{theorem}

In terms of \(n\) the bound is \(2^{\Omega(n/\log^2n)}\). We call it
exponential in the usual sense of proof complexity: the formula
\eqref{eq:bit-php} has size \(L=\Theta(n^3\log n)\), and the bound is
\(2^{L^{1/3-o(1)}}\). No bound of the form \(2^{\Omega(L)}\) or
\(2^{\Omega(n)}\) is claimed. The constant
\(32768=2^{15}\) and the threshold \(32\) are what the formal proof
verifies; no attempt was made to optimize them. Version 1 of this preprint
stated only the following qualitative form, which was formalized first and
which the main theorem implies for large \(\ell\).

\begin{corollary}[Superpolynomial form]\label{cor:main-superpolynomial}
\lean{claims/BitPHPSuperpolynomial.lean}
For every fixed real \(K>0\), there is \(\ell_0(K)\) such that for every
\(\ell\ge\ell_0(K)\), every DAG-like \(\Res\) refutation of
\(\BPHP_{2^\ell}^{2^\ell+1}\) has more than \(2^{K\ell}=n^K\) nodes.
\end{corollary}

The input length in \eqref{eq:bit-php} is polynomial in \(n\), so the bound
is also exponential in a fixed power of the input length, and it bounds the
total proof-description size from below.

\subsection{Prior results and the scope of the claim}
Efremenko, Garl\'ik, and Itsykson \cite{EGI25} prove a
\(2^{\Omega(n^{1/3}/\log n)}\) size bound for regular \(\Res\) refutations
of bit PHP. Their unrestricted DAG-like result is an affine-width bound.
Bhattacharya, Chattopadhyay, and Dvo\v{r}\'ak \cite{BCD24} exhibit formulas
that are exponentially hard for bottom-regular \(\Res\), although they
have short proofs in general ordinary resolution. This separation
demonstrates the importance of removing regularity.

Bhattacharya and Chattopadhyay \cite{BC25}, Byramji and Impagliazzo
\cite{BI25}, and the merged STOC 2026 paper \cite{BCBI26} obtain
exponential bounds in regimes allowing nearly quadratic proof depth.
Their bounds retain a proof-depth restriction. Itsykson, Podolskii,
and Shekhovtsov \cite{IPS26} explicitly distinguish this progress from
the general superpolynomial size problem; their unrestricted size
consequence is quadratic.

Alekseev and Gaevoy \cite{AG26} study constrained bit PHP. Their
polynomial-depth bounds for general \(\Res\) are conditional; their
unconditional results retain either a reversible-system restriction or
a near-quadratic depth bound. They therefore do not give the unrestricted
usual-bit-PHP size conclusion considered here.

\paragraph{Resolution over linear equations.}
The family \(\mathrm{Res}(\mathrm{lin}_R)\) studies resolution over linear
equations in a ring \(R\). Raz and Tzameret \cite{RT08} introduced the
integer formulation; Itsykson and Sokolov \cite{IS14,IS20} studied the
\(\F\) version in the tree-like setting. Part and Tzameret
\cite{PT21} proved the first superpolynomial lower bounds for DAG-like
\(\mathrm{Res}(\mathrm{lin}_R)\): the subset-sum principle requires
exponential size over the rationals. Their technique uses large
characteristic and does not apply over \(\F\). Khaniki \cite{Kha22} proved
almost quadratic lower bounds for DAG-like resolution over low-degree
polynomial equations over finite fields, which includes resolution over
linear equations, for mod-\(q\) Tseitin formulas and random \(k\)-CNFs.
To our knowledge these were the first superlinear DAG-like bounds for such
systems over finite fields; \cite{EGI25} notes that the rules considered
there differ from the formulation of \(\Res\) used here.

\paragraph{Other work on \(\Res\).}
Gryaznov, Pudl\'ak, and Talebanfard \cite{GPT22} introduced regular \(\Res\)
and related it to read-once linear branching programs. Alekseev and Itsykson
\cite{AI25} lift resolution depth to regular and bounded-depth \(\Res\) size,
and Efremenko and Itsykson \cite{EI25,EI26} sharpen the closure technique
of \cite{EGI25} and obtain strong bounds for bounded-depth \(\Res\). All of
these retain a regularity or depth restriction.

\paragraph{Polynomial calculus and the pigeonhole principle.}
The degree lower bound at the base of our argument is Razborov's theorem
\cite{Raz98} that polynomial calculus refutations of the pigeonhole principle
need degree \(n/2+1\) over every field, which also describes the low-degree
consequences of the pigeonhole axioms. Related degree bounds and methods are
due to Impagliazzo, Pudl\'ak, and Sgall \cite{IPS99}, Alekhnovich and Razborov
\cite{AR03}, and Mik\v{s}a and Nordstr\"om \cite{MN15}. Section~\ref{sec:moments}
reproves the part of Razborov's theorem that we need, over \(\F\), by a
different route through chessboard complexes; we do not claim this degree
range as new (see the discussion before Theorem~\ref{thm:moments}).
The extension variables are those of Buss, Impagliazzo, Kraj\'i\v{c}ek,
Pudl\'ak, Razborov, and Sgall \cite{BIKPRS}.

\paragraph{Scope.}
The theorem here removes both regularity and proof-depth restrictions.
Targeted searches of the primary literature on 15 and 17 September 2026 found
no preceding theorem with that scope; this is not an exhaustive certification
of novelty. The theorem concerns \(\Res\) only. It does not address unary
PHP in \(\Res\), and the \(\mathrm{AC}^0[p]\)-Frege lower-bound problem
remains open; Section~\ref{sec:generic} says what does and does not follow.

\subsection{Revision history}
Version 1 (15 September 2026) stated Corollary~\ref{cor:main-superpolynomial}.
Revision 1 adds the exponential bound of Theorem~\ref{thm:main}
(Section~\ref{sec:main-proof}) and the generic sufficient condition of
Section~\ref{sec:generic}, both formalized, a proof overview
(Section~\ref{sec:overview}), and a fuller account of related work. The
development history formerly in the last section was condensed. No error in
the mathematics of version 1 was found or corrected; the registry slot count
of Section~\ref{sec:clause-simulation} was renamed from \(N\) to
\(N_{\mathrm{reg}}\) to avoid a clash with the number of unary columns.

Revision 2 (19 September 2026) corrects the historical attribution of
resolution over linear equations, makes the common-kernel and
bounded-cofactor steps more explicit, and improves the parameter tables.
Appendix~\ref{sec:verification} adds a formal-verification map and complete
reproduction instructions within the paper. The theorem statements,
mathematical arguments, and Lean dependencies are unchanged from revision 1;
this revision makes no new mathematical claim.

\subsection{Organization}
Section~\ref{sec:overview} gives an overview of the proof, fixes terminology
and parameters, and works through one simulated inference.
Sections~\ref{sec:algebra}--\ref{sec:clause-simulation} prove the components
in full, and Section~\ref{sec:main-proof} closes the parameters.
Section~\ref{sec:generic} states the generic criterion; it can be skipped.
Appendix~\ref{sec:topology} proves the homological input behind
the matching extension. A companion online research notebook records the
development of the argument and supporting experiments \cite{Repo};
Section~\ref{sec:framework} describes these supplementary materials.
The proofs presented here can be read independently of that notebook:
none of its lemmas is required as an unexpanded black box in this preprint.

Every [Lean] link points to a source file in the published revision
\href{https://github.com/kbr-/math-research/tree/8904bf09a5376f48a00d1c25d079bf0c6441502a}{\texttt{8904bf09}}, rather than a moving
branch. The linked files identify their exported statements and dependencies.
Appendix~\ref{sec:verification} maps the main proof steps to exact Lean
declarations and gives the checkout, setup, and verification commands.
The mathematical notation here is independent of Lean syntax.
